\subsection{Messverfahren}
Zuerst wird durch Messen der Höhe und Länge der schiefen Ebene mit einem Meter-Lineal die Größen zur Berechnung des Neigungswinkels erhoben. Außerdem wurde die Massen und Durchmesser der Zylinder mit dem Messschieber erfasst.\\
Anschließend wurden zur Überzeugung, dass das Abrollen an der schiefen Ebene sich um eine beschleunigte Bewegung handelt, die Messsensoren auf Positionen \(P_i=\)\qtylist[list-units=bracket]{9;25;49;81}{\cm} gestellt. Wenn \cref{eq:s(t^2)} gilt, also \(t=\sqrt{\frac{2s}{a}}\), und für \(s=k^2\) mit \(k\in(3,5,7,9)\) eingesetzt wurde, sollte sich also eine lineare Zunahme der Zeit mit \(t=\sqrt{\frac{2}{a}}k\) ergeben bzw. also eine konstante Zeitdifferenz zwischen den \(t_i(P_i)\) auftreten. Dies konnte bestätigt werden, siehe Spalte \emph{Differenz} in \cref{table:aufgabe1mittelwerte}. Die erhobenen Messwerte wurden in \cref{sec:aufgabe1} bzw. \cref{fig:diagramm1} ausgewertet.

Als letzter Schritt wurde eine Zeitmessung für unterschiedliche Lichtschrankenpositionen auf der Auslauffläche durchgeführt.

Alle Zeitmessungen erfolgten durch elektronische Stoppuhren, die automatisch bei Anheben der Startschranke gestartet und bei Auslösung der Lichtschranken gestoppt wurden.
\vfill
\subsection{Messprotokoll}
\begin{figure}[htbp]
    \centering
    \includegraphics[page=1,width=\textwidth]{V15.pdf}
\end{figure}
\newpage

\begin{table}[t]
    \caption{Tabelle 1 des Messprotokolls: Zeiten an der schiefen Ebene}
    \includegraphics[page=1,width=\textwidth]{V15_1.pdf}
\end{table}

\begin{table}[htbp]
    \caption{Tabelle 2 des Messprotokolls: Zeiten an der Auslauffläche}
    \includegraphics[page=1,width=\textwidth]{V15_2.pdf}
\end{table}
\clearpage
